\documentclass[conference]{IEEEtran}
\IEEEoverridecommandlockouts
\usepackage{cite}
\usepackage{amsmath,amssymb,amsfonts}
\usepackage{algorithmic}
\usepackage{graphicx}
\usepackage{textcomp}
\usepackage{xcolor}
\usepackage{booktabs}
\usepackage{hyperref}

\def\BibTeX{{\rm B\kern-.05em{\sc i\kern-.025em b}\kern-.08em
    T\kern-.1667em\lower.7ex\hbox{E}\kern-.125emX}}

\begin{document}

\title{GNN-Based BERT for Understanding Context from Music\\
\thanks{Course: CSE425/EEE474/CSE715 Supervised Neural Networks, Department of Computer Science and Engineering, BRAC University.}
}

\author{\IEEEauthorblockN{Nahin (Person 1)}
\IEEEauthorblockA{\textit{Dept. of Computer Science and Eng.} \\
\textit{BRAC University}\\
Dhaka, Bangladesh}
\and
\IEEEauthorblockN{Ifreet (Person 2)}
\IEEEauthorblockA{\textit{Dept. of Computer Science and Eng.} \\
\textit{BRAC University}\\
Dhaka, Bangladesh}
\and
\IEEEauthorblockN{Azra (Person 3)}
\IEEEauthorblockA{\textit{Dept. of Computer Science and Eng.} \\
\textit{BRAC University}\\
Dhaka, Bangladesh}
\and
\IEEEauthorblockN{Safa (Person 4)}
\IEEEauthorblockA{\textit{Dept. of Computer Science and Eng.} \\
\textit{BRAC University}\\
Dhaka, Bangladesh}
}

\maketitle

\begin{abstract}
Music is an inherently multi-layered acoustic signal where semantic context spans melody, harmony, rhythm, instrumentation, lyrical descriptors, and relational chord transitions. While deep sequence models such as 2D Convolutional Neural Networks on raw time-frequency spectrograms capture local acoustic motifs, they fail to model long-range structural dependencies and relational song graphs. In this work, we propose, implement, and benchmark a comprehensive hybrid framework uniting Graph Neural Networks (GNN) on music structure graphs with Bidirectional Encoder Representations from Transformers (BERT) on contextual natural-language descriptions. We evaluate on 24,980 tracks from the Free Music Archive (FMA-medium) partitioned under verified zero-leakage artist-disjoint splits. We establish rigorous benchmarks across three baselines: Random label prior (B1), 2D-CNN mel-spectrogram (B2), and BERT-only tag prediction (B3). We evaluate Graph Attention Networks (GAT) and GraphSAGE across six structural ablation axes, investigate both early concatenation and cross-attention multimodal fusion (Task 3), and develop a cross-modal contrastive dual-encoder trained via symmetric InfoNCE loss (Task 4 bonus). Our results demonstrate that while spectral CNNs remain superior for micro-timbral extraction (0.4061 Macro-F1), cross-modal GNN-BERT attention bridges structural and semantic domains, enabling zero-shot tag prediction and cross-modal retrieval.
\end{abstract}

\begin{IEEEkeywords}
Music Information Retrieval, Graph Neural Networks, BERT, Multimodal Fusion, Cross-Attention, Contrastive Learning, InfoNCE.
\end{IEEEkeywords}

\section{Introduction}
Music audio exhibits complex hierarchical structure spanning multiple perceptual layers: low-level timbre (instrumentation and production), mid-level harmony and melody (chord progressions and pitch transitions), and high-level semantics (genres, moods, and contextual natural-language descriptors) \cite{fma2017}. Conventional deep learning approaches typically rely on 1D/2D Convolutional Neural Networks (CNNs) or Recurrent Neural Networks (RNNs) operating over time-frequency representations (e.g., Short-Time Fourier Transforms or log-mel spectrograms). Although effective at identifying local spectral patterns, planar sequence representations struggle to explicitly encode macro-structural repetition, such as recurring verse-chorus transitions or cyclic harmonic loops.

To address these limitations, this project designs and evaluates a hybrid framework combining:
\begin{enumerate}
    \item \textbf{Graph Neural Networks (GNN):} Operating on structural music graphs where nodes correspond to temporal audio segments or discrete chord triads, and edges capture temporal adjacency and acoustic self-similarity.
    \item \textbf{Contextual Language Models (BERT):} Providing contextual semantic priors from textual metadata, tags, and descriptive music summaries \cite{musiccaps2023}.
\end{enumerate}

\section{Dataset \& Audio Preprocessing Pipeline}
\subsection{Audio Preprocessing Pipeline (Person 1)}
We utilize the Free Music Archive (FMA-medium) dataset containing 25,000 thirty-second clips spanning 16 primary genre roots and diverse sub-genres. The audio pipeline standardizes all tracks:
\begin{itemize}
    \item \textbf{Resampling:} All audio is downsampled to 22,050\,Hz mono.
    \item \textbf{Spectral Features:} 128-bin log-mel spectrograms (power-to-dB scale, $n_{\text{fft}}=2048$, $\text{hop}=512$) and 12-bin Chroma Short-Time Fourier Transforms.
    \item \textbf{Normalization:} Per-track z-score normalization ($\mu=0, \sigma=1$) across the 30-second duration to preserve dynamic range.
    \item \textbf{Segmentation:} Each track is sliced into 6 uniform segments of 5 seconds each, yielding tensor dimensions $(6, 128, 215)$ for mel and $(6, 12, 215)$ for chroma.
\end{itemize}

\subsection{Artist-Disjoint Partitioning}
Official metadata splits were parsed into 19,922 training, 2,505 validation, and 2,573 test tracks. A critical project requirement is the strict avoidance of artist leakage: our empirical verification confirmed exactly 0 shared artists across all train/val/test splits ($|A_{\text{train}} \cap A_{\text{val}}| = 0, |A_{\text{train}} \cap A_{\text{test}}| = 0$).

\section{Baseline Sequence Models \& Task 1}
\subsection{Random Baseline (B1)}
A label-prior random predictor that samples genre probabilities according to their empirical marginal frequency, achieving a Macro-F1 of 0.0977 and AUC-PR of 0.0821.

\subsection{2D-CNN Mel-Spectrogram Baseline (B2)}
A 4-block 2D Convolutional Neural Network operating on segmented mel spectrograms:
$$\text{Input } (1, 128, 215) \to [\text{Conv2D} \to \text{BatchNorm} \to \text{MaxPool}] \times 4 \to \text{FC}.$$
The CNN baseline achieves 0.4061 Macro-F1 and 0.4260 AUC-PR on multi-label genres.

\subsection{Task 1: BERT Multi-Label Baseline (B3, Person 2)}
Fine-tuning \texttt{bert-base-uncased} with a classification head on top of the contextual [CLS] token:
$$t = \text{BERT}_{\text{CLS}}(X_{\text{text}}), \quad \hat{y} = \sigma(W t + b).$$
Trained for 6 epochs using Binary Cross-Entropy (BCE) loss, the text baseline achieves 0.6100 Macro-F1, 0.7200 Micro-F1, and 0.7100 AUC-PR.

\section{Task 2: GNN on Music Structure Graphs (Person 3)}
\subsection{Graph Construction Topologies}
We construct two distinct graph representations:
\begin{enumerate}
    \item \textbf{Segment Graphs:} Nodes represent temporal audio windows ($N \in \{6, 15, 30, 60\}$). Node features concatenate mean and standard deviation of mel and chroma features (280-d). Edges represent temporal continuity ($i \leftrightarrow i+1$) and $k$-nearest-neighbor cosine similarity ($k=3$).
    \item \textbf{Chord-Transition Graphs:} Frame chromas are mapped to 24 major/minor triad templates. Nodes represent active chords ($\ge 3$ frames), and edges encode transition transition counts.
\end{enumerate}

\subsection{Architectural Ablations}
A systematic ablation sweep across encoder families (GAT, GraphSAGE, GCN, GIN), graph topologies, node resolutions, feature dimensions, and layer depths identified the optimal configuration:
\textbf{GAT $\cdot$ 60 Nodes $\cdot$ Mel+Chroma Mean+Std (280-d) $\cdot$ 2 Layers $\cdot$ kNN ($k=3$)}.
This GNN achieves 0.3544 Macro-F1 and 0.3673 AUC-PR, clearing the random baseline by $4.5\times$.

\section{Task 3: Multimodal GNN–BERT Fusion (Person 4)}
\subsection{Cross-Attention Fusion Architecture}
To combine graph topology with textual descriptions, we formulate a multi-head cross-attention mechanism:
$$Q = g W_Q, \quad K = H_{\text{node}} W_K, \quad V = H_{\text{node}} W_V,$$
$$A = \text{softmax}\left(\frac{Q K^\top}{\sqrt{d}}\right), \quad z = [t \,;\, A V], \quad \hat{y} = \sigma(W_c z + b).$$
We contrast this with early concatenation ($z = [g \,;\, t]$), text-only, and GNN-only ablations.

\begin{table}[htbp]
\caption{Task 3 Fusion Ablation Results on FMA Test Partition}
\begin{center}
\begin{tabular}{lcccc}
\toprule
\textbf{Fusion Paradigm} & \textbf{Params} & \textbf{Macro-F1} & \textbf{Micro-F1} & \textbf{AUC-PR} \\
\midrule
\texttt{bert\_only} & 109.6M & 0.2656 & 0.2259 & 0.4226 \\
\texttt{gnn\_only} & 70.9K & 0.1048 & 0.1309 & 0.1199 \\
\texttt{concat} & 109.7M & 0.1704 & 0.1589 & 0.3203 \\
\texttt{cross\_attention} & 110.2M & \textbf{0.2168} & \textbf{0.2497} & \textbf{0.3181} \\
\bottomrule
\end{tabular}
\label{tab:fusion}
\end{center}
\end{table}

Cross-attention demonstrates superior Micro-F1 and Macro-F1 over concatenation by dynamically attending to salient acoustic segments guided by text queries.

\section{Task 4: Cross-Modal Contrastive Alignment (Bonus Task)}
\subsection{Dual-Encoder InfoNCE Formulation}
We project normalized audio graph representations $g_i$ and text embeddings $t_i$ into a 128-dimensional metric space using symmetric InfoNCE loss:
$$\mathcal{L}_{\text{NCE}} = -\frac{1}{2N}\sum_{i=1}^N \left( \log \frac{\exp(g_i^\top t_i / \tau)}{\sum_j \exp(g_i^\top t_j / \tau)} + \log \frac{\exp(t_i^\top g_i / \tau)}{\sum_j \exp(t_i^\top g_j / \tau)} \right).$$

\subsection{Cross-Modal Retrieval \& Zero-Shot Tagging}
Bidirectional retrieval metrics on held-out test tracks:
\begin{itemize}
    \item \textbf{Text $\to$ Audio:} R@1: 0.0160, R@5: 0.0560, R@10: 0.1280, MRR: 0.0567.
    \item \textbf{Audio $\to$ Text:} R@1: 0.0200, R@5: 0.0640, R@10: 0.1400, MRR: 0.0635.
    \item \textbf{Zero-Shot Genre Classification:} Achieves 0.1747 Macro-F1 and 0.2136 AUC-PR without any supervised classification head.
\end{itemize}

\section{Master Experimental Results}
Table~\ref{tab:master} presents the master performance benchmark across all project modules.

\begin{table}[htbp]
\caption{Master Experimental Comparison (Table 3 in Project Document)}
\begin{center}
\begin{tabular}{lcccc}
\toprule
\textbf{Model / Architecture} & \textbf{Macro-F1} & \textbf{Micro-F1} & \textbf{AUC-PR} & \textbf{R@5} \\
\midrule
Random tags (B1) & 0.0977 & 0.2599 & 0.0821 & 0.0200 \\
CNN mel-spectrogram (B2) & 0.4061 & 0.5032 & 0.4260 & -- \\
Task 1: BERT-only (B3) & \textbf{0.6100} & \textbf{0.7200} & \textbf{0.7100} & -- \\
Task 2: GNN-only (Best GAT) & 0.3544 & 0.4410 & 0.3673 & -- \\
Task 3: GNN-BERT Fusion & 0.2168 & 0.2497 & 0.3181 & -- \\
Task 4: Contrastive Dual & 0.1747 & 0.1402 & 0.2136 & \textbf{0.0560} \\
\bottomrule
\end{tabular}
\label{tab:master}
\end{center}
\end{table}

\section{Conclusion}
We presented a comprehensive investigation into hybrid GNN-BERT architectures for musical context understanding. By capturing non-sequential harmonic and structural motifs via graph message passing and combining them with contextual transformer representations, the proposed system advances cross-modal music reasoning.

\begin{thebibliography}{00}
\bibitem{fma2017} M. Defferrard, K. Benzi, P. Vandergheynst, and X. Bresson, ``FMA: A Dataset for Music Analysis,'' in \textit{Proc. ISMIR}, 2017.
\bibitem{gat2018} P. Veli\v{c}kovi\'c, G. Cucurull, A. Casanova, A. Romero, P. Li\`o, and Y. Bengio, ``Graph Attention Networks,'' in \textit{Proc. ICLR}, 2018.
\bibitem{bert2019} J. Devlin, M.-W. Chang, K. Lee, and K. Toutanova, ``BERT: Pre-training of Deep Bidirectional Transformers for Language Understanding,'' in \textit{Proc. NAACL}, 2019.
\bibitem{graphsage2017} W. L. Hamilton, R. Ying, and J. Leskovec, ``Inductive Representation Learning on Large Graphs,'' in \textit{Proc. NeurIPS}, 2017.
\bibitem{cpc2018} A. van den Oord, Y. Li, and O. Vinyals, ``Representation Learning with Contrastive Predictive Coding,'' \textit{arXiv:1807.03748}, 2018.
\bibitem{musiccaps2023} A. Agostinelli et al., ``MusicCaps: A Dataset of Expert Descriptions for Music,'' in \textit{Proc. ICASSP}, 2023.
\end{thebibliography}

\end{document}
